\scanpage{109}
Les sous fibrés vectoriels $B$ de $A$, \inserted{contt $1$,} de rang $3$, correspondent par orthogonalité (pour $\varphi$ ou $\psi$, c'est kif kif) aux sous fibrés vectoriels de rang $1$ $D$ de $A^{\unsure{s|\leftrightarrow|\natural}}$.

D'autre part, pour un tel $B$, considérons
\[
\det_B=\det\mid B\ ,\quad\text{forme quadratique sur } B.
\]
elle admet un discriminant divisé
\[
\deleted{\unsure{\delta}}\operatorname{disc}'(\det_B)\quad\text{tels que}\quad2\operatorname{disc}'(\det_B)=\operatorname{disc}\psi_B
\]
où $\psi_B$ est la restriction de $\psi$ à $B$.

Ainsi
\[
\deleted{\unsure{2}}\operatorname{disc}\psi_B=\operatorname{disc}\varphi_B
=2\underbrace{\operatorname{disc}'(\det_B)}_{\substack{\parallel\,\text{déf}\\\delta(B)}}\unsure{=}
\]

\uline{Proposition} Les \inserted{quatre} conditions suivantes sont équivalentes

(i) $D$ est \uline{isotrope} i.e. $\det\mid D=0$

(ii) $\delta(B)=0$\qquad (\uline{NB} $\delta(B)\in\deleted{\unsure{(\det B)^{\otimes2}}}\det(B)^{\otimes(-2)}$).

\deleted{\unsure{(iii)}} (iii) $B$ est une sous-algèbre de $\unsure{A|M}$

% slanted
\margin{\deleted{Si $k$ un corps \unsure{sinon|il suffit} on sait pas si ça implique (iii$'$) \unsure{et}} $\longrightarrow$}

(iii$'$) \deleted{\unsure{(iii$'$)}} $B$ est une ``sous-alg. de Borel'' de $A$

Ces \deleted{\unsure{(iv)}} conditions impliquent les deux suivantes, et si $2.1_k$ est inv, \deleted{elles} ces dernières sont équivalentes aux précédentes

(iv) $D\subset B$

(\unsure{iv$'$|iii$'$|ii$'$}) \deleted{\unsure{$\operatorname{disc}=0$|$\operatorname{disc}$}} $\operatorname{disc}(\varphi_B)=0$\qquad ($\Longleftrightarrow\operatorname{disc}\psi_B=0$)
% archive sheet number: 108
